% P17-D5: Rubinstein Lecture 9, Problem 2, and proved median consequences.
\chapter{Single-Peaked Preferences and Majority Choice}
\label{sec:p17-single-peaked-majority}

Majority comparison becomes a consistent social ranking on an important
restricted preference domain. The alternatives share one specified order,
and each voter prefers movement toward one personal peak on either side
of it. With an odd finite electorate and strict individual rankings, the
entire majority relation is transitive. The middle peak also defeats every
other alternative in a pairwise vote. These are separate conclusions, and
both receive full proofs below.

The source problem is Rubinstein's Lecture 9, Problem 2, in the author
revision dated 11 January 2015. Its alternative set is the whole interval
\([0,1]\), and it inherits Lecture 9's strict individual-order domain
\cite[pp.~123, 128]{Rubinstein2015LectureNotes}. We retain that regime.
The argument also gives the explicitly stated results for other ordered
sets. The median, finite procedure and parity refinements are deductions
supplied here; the source poses the odd-electorate aggregation result
as a problem, rather than printing its solution.

\section{The common order and the restricted profile domain}
\label{sec:p17-single-peaked-domain}

\begin{hbdefinition}{def:p17-single-peaked-order}{A peak relative to a specified axis}
Let \(X\) be a nonempty set with a specified linear order \(\leq\),
called the \emph{axis}; write \(<\) for its strict part. This is an
order of alternatives, not a voter's preference relation. In the source
regime, \(X=[0,1]\) with the ordinary real order.

A strict ordering \(\succ_i\in\mathcal L(X)\), as defined in
Definition~\ref{def:p17-social-profile}, is \emph{single-peaked relative
to the axis} if there is \(p_i\in X\) such that, for all \(b,c\in X\),
\[
 (p_i\geq b>c\ \hbox{or}\ c>b\geq p_i)
                 \quad\Longrightarrow\quad b\succ_i c.
\]
This is the source condition, with \(p_i\) denoting its \(a_i^*\).
It requires strict improvement toward the peak on either side, including
comparison with the peak itself. No continuity, numerical utility,
symmetry about the peak or distance-based comparison across opposite
sides is assumed. The peak must belong to \(X\); its existence is part
of the definition even when \(X\) is infinite.

Write \(\mathcal L_{\mathrm{SP}}(X,\leq)\) for this subset of
\(\mathcal L(X)\). For \(N=\{1,\ldots,n\}\), \(n\geq1\), the
restricted profile domain is
\[
 \mathcal D_{\mathrm{SP}}\defeq
                   \mathcal L_{\mathrm{SP}}(X,\leq)^N.
\]
Every voter uses this same fixed axis. A different axis for each voter
does not satisfy this domain requirement.
\index{single-peaked preference}\index{preference!common axis}
\end{hbdefinition}

\begin{hblemma}{thm:p17-single-peaked-structure}{proved=annex:p17.single-peaked-structure}{Unique peaks and the four rankings of three ordered alternatives}
Every single-peaked strict ordering has a unique peak \(p_i\), and
\(p_i\succ_i x\) for every \(x\ne p_i\). For any three alternatives
\(a<b<c\), their middle member cannot be ranked below both endpoints.
Consequently the restriction of an individual ordering to this triple
has exactly one of the following four forms:
\[
 a\succ_i b\succ_i c,\qquad
 b\succ_i a\succ_i c,\qquad
 b\succ_i c\succ_i a,\qquad
 c\succ_i b\succ_i a.
\]
Conversely, a strict ordering with a greatest preferred alternative \(p\)
is single-peaked if no triple \(a<b<c\) ranks \(b\) below both
endpoints. The premise that a greatest preferred alternative exists is
automatic on a finite nonempty \(X\), but is retained on infinite sets.
\end{hblemma}

\begin{hbdefinition}{def:p17-single-peaked-vote}{Majority, ties and Condorcet winners}
At \(P\in\mathcal D_{\mathrm{SP}}\), use the same vote counts and
majority convention as Definition~\ref{def:p17-pairwise-majority}:
\[
 V_P(x,y)=\#\{i\in N:x\succ_i y\}\qquad(x\ne y).
\]
For distinct \(x,y\), the relation \(x\succ^M_P y\) means
\(V_P(x,y)>V_P(y,x)\); no \(x\succ^M_Px\) holds.
Its weak version \(x\succeq^M_Py\)
uses \(\geq\) on distinct pairs and contains every diagonal pair.
Equal counts give \(x\sim^M_Py\). These are social ties; the
individual inputs remain strict, so
\(V_P(x,y)+V_P(y,x)=n\).

A \emph{strict Condorcet winner} is an alternative \(w\in X\)
such that \(w\succ^M_Px\) for every \(x\ne w\).
A \emph{weak Condorcet winner} satisfies \(w\succeq^M_Px\) for
every \(x\in X\), equivalently no alternative strictly defeats it.
These definitions concern all pairwise votes; neither definition
presupposes that the whole majority relation is transitive.
On a singleton alternative set its unique member satisfies both
definitions. A strict Condorcet winner, if one exists, is unique.
\index{Condorcet winner}\index{majority!social ties}
\end{hbdefinition}

\section{The whole majority ranking and its greatest alternative}
\label{sec:p17-single-peaked-theorems}

\begin{hbtheorem}{thm:p17-single-peaked-majority}{proved=annex:p17.single-peaked-majority}{Odd majority is a transitive social ordering}
Let \(X\) be a nonempty linearly ordered set, let \(n=2k+1\) with
\(k\geq0\), and let \(P\in\mathcal D_{\mathrm{SP}}\) have strict
individual inputs. Then \(\succ^M_P\) is a strict ordering: every two
distinct alternatives are compared in exactly one direction, and it is
transitive. Its reflexive closure \(\succeq^M_P\) is complete and
transitive, with no tie between distinct alternatives. Therefore
\[
 F_{\mathrm{SP}}:\mathcal D_{\mathrm{SP}}\longrightarrow\mathcal W(X),
 \qquad F_{\mathrm{SP}}(P)=\succeq^M_P,
\]
is a social welfare function on the restricted domain. In particular
this proves the source's conclusion for every odd positive \(n\) on
the whole interval \([0,1]\).
\end{hbtheorem}

\begin{hbdefinition}{def:p17-median-peaks}{The median of a finite list of peaks}
Sort the \(n\) peaks in axis order, retaining repeated occurrences:
\[
 p_{(1)}\leq p_{(2)}\leq\cdots\leq p_{(n)}.
\]
Parenthesized subscripts refer to positions in this sorted list, not
voter identities. A finite list in a linearly ordered set can be sorted
without an order-completeness or topological assumption on that set.
For \(n=2k+1\), its \emph{median peak} is \(m=p_{(k+1)}\).
For \(n=2k\), \(k\geq1\), its two central peaks are
\(l=p_{(k)}\) and \(u=p_{(k+1)}\).
For \(l\leq u\) in \(X\), let
\([l,u]_X=\{x\in X:l\leq x\leq u\}\).
\index{median voter}\index{median!peak}
\end{hbdefinition}

\begin{hbtheorem}{thm:p17-median-peak-winner}{proved=annex:p17.median-peak-winner}{The median peak wins; the majority ordering is single-peaked}
Let \(n=2k+1\) with \(k\geq0\), and let \(P\) be a profile of
strict single-peaked orderings on the same nonempty linearly ordered
\(X\). Then the median peak \(m\) is the unique strict
Condorcet winner. It is also the unique weak Condorcet winner and the
greatest alternative of the majority ordering proved above.
For every \(x<m\), at least \(k+1\) voters strictly prefer \(m\)
to \(x\); for every \(x>m\), at least \(k+1\) voters strictly
prefer \(m\) to \(x\).

Moreover, \(\succ^M_P\) is itself single-peaked on this same axis
with peak \(m\). Repeated individual peaks, endpoint peaks in
\([0,1]\), and the one-voter case are included. The median may be
the peak of more than one voter; no fixed voter is selected by the
order-statistic notation.
\end{hbtheorem}

\begin{hbtheorem}{thm:p17-even-median-interval}{proved=annex:p17.even-median-interval}{Even electorates: the exact weak-winner interval}
Let \(n=2k\), \(k\geq1\), and keep strict single-peaked inputs
on the same nonempty linearly ordered \(X\). The set of weak Condorcet
winners is exactly \([l,u]_X\), where \(l,u\) are the two central
peaks. There is a strict Condorcet winner exactly when \(l=u\);
in that case it is their common value. This conclusion does not assert
transitivity of the weak majority relation for an even electorate.
\end{hbtheorem}

\begin{hbexample}{def:p17-single-peaked-profile}{Five strict rankings and the complete majority result}
Take the five-alternative axis
\(a=0<b=1/4<c=1/2<d=3/4<e=1\). In the following table each
row lists a voter's alternatives from most to least preferred:
\[
 \begin{array}{c|ccccc|c}
 \text{voter}&1\text{st}&2\text{nd}&3\text{rd}&4\text{th}&5\text{th}&\text{peak}\\\hline
 1&a&b&c&d&e&a\\
 2&b&a&c&d&e&b\\
 3&c&b&d&e&a&c\\
 4&d&c&b&a&e&d\\
 5&e&d&c&b&a&e
 \end{array}
\]
Each ranking strictly improves toward its listed peak on each side.
The matrix below records \(V_P(x,y)\), with \(x\) the row and
\(y\) the column; diagonal entries are unused:
\[
 \begin{array}{c|ccccc}
      &a&b&c&d&e\\\hline
 a&-&1&2&2&3\\
 b&4&-&2&3&4\\
 c&3&3&-&3&4\\
 d&3&2&2&-&4\\
 e&2&1&1&1&-
 \end{array}
\]
Thus the entire strict majority ordering is
\(c\succ^M_P b\succ^M_P d\succ^M_P a\succ^M_P e\).
The sorted peaks are \(a,b,c,d,e\), so the median \(c\) defeats
every rival. The majority ordering differs from voter 3's ordering:
the majority prefers \(a\) to \(e\), whereas voter 3 prefers
\(e\) to \(a\). A median-peak winner theorem does not make the
median-peak voter's entire ranking socially decisive. The annex verifies
all pair counts in this example with the other boundary calculations.
\end{hbexample}

\section{A finite procedure and its limits}
\label{sec:p17-single-peaked-procedure}

\begin{hbmethod}{def:p17-single-peaked-finite-check}{Checking a profile against a supplied finite axis}
The input is a nonempty list of distinct alternatives in axis order,
\(a_1<\cdots<a_q\), and \(n\) preference lists, each claiming to
be a permutation of those alternatives from best to worst. The axis
is supplied; finding an axis on which the profile becomes single-peaked
is a different problem.
\begin{enumerate}
\item Verify that each preference list contains every index from one
      through \(q\) exactly once. Build the rank array \(r_i(j)\),
      with rank one best, and let \(h_i\) be the axis index of its
      top alternative.
\item For every voter check
      \[
       r_i(j)>r_i(j+1)\quad(1\leq j<h_i),\qquad
       r_i(j)<r_i(j+1)\quad(h_i\leq j<q).
      \]
      A failed comparison supplies an adjacent same-side pair violating
      the peak condition. Return that voter, peak and pair as a witness.
\item After all checks pass, count the occurrences of each peak index.
      Cumulative counts locate the middle peak for odd \(n\), or the
      two central peaks for even \(n\). Return the strict winner in
      the odd case and the interval of weak winners in the even case.
      In the even case, also report a strict winner exactly when the
      two central peaks coincide.
\item If the full majority relation is requested, count votes for every
      unordered pair by comparing its two rank entries for every voter.
      For odd \(n\), this is the complete transitive ranking. One can
      list it in order by each alternative's number of strict pairwise
      victories, from \(q-1\) down to zero.
\end{enumerate}
This procedure acts on a finite encoding. A finite list of observations
on \([0,1]\) is not a certificate that a voter's entire interval
ordering satisfies the universally quantified source condition.
\end{hbmethod}

\begin{hbtheorem}{thm:p17-single-peaked-algorithm}{proved=annex:p17.single-peaked-algorithm}{Correctness and finite work of the supplied-axis procedure}
For \(q,n\geq1\), the procedure rejects invalid preference-list
encodings and otherwise accepts exactly the profiles that are
single-peaked on the supplied axis. Its peak and winner outputs are
correct. For odd \(n\), its pairwise matrix and victory-count ordering
give the whole strict majority ordering.

With alternatives encoded by their axis indices, random-access rank
arrays, and unit-cost operations on the indices and vote counts,
validation and the single-peaked checks take \(O(nq)\) operations.
After the peaks are available, finding the central peak or peaks by
their histogram takes \(O(n+q)\) operations. Computing every pairwise
count uses exactly \(nq(q-1)/2\) individual rank comparisons, and
the complete ranking can be returned in \(O(nq^2)\) operations.
These are bounds in the stated finite representation, not a decision
procedure for arbitrary descriptions of preferences on an interval.
\end{hbtheorem}

\begin{hbtheorem}{thm:p17-single-peaked-boundaries}{proved=annex:p17.single-peaked-boundaries}{The worked profile and two failed hypotheses}
The profile and count matrix in Example~\ref{def:p17-single-peaked-profile}
have the stated properties. Two further finite examples locate the domain
and parity requirements:
\begin{enumerate}
\item For \(x<y<z\), take three voters with rankings
      \(x\succ y\succ z\), \(y\succ z\succ x\), and
      \(z\succ x\succ y\). Their majority has the strict cycle
      \(x\succ^M y\succ^M z\succ^M x\), by two votes to one
      in each comparison. The third ranking violates single-peakedness
      on this axis because it puts the middle alternative \(y\) last.
\item On the same axis take four voters: two have
      \(x\succ y\succ z\), one has \(y\succ z\succ x\),
      and one has \(z\succ y\succ x\). All four rankings are
      single-peaked. The majority ties \(x\) with \(y\) and
      ties \(x\) with \(z\), while \(y\succ^M z\) by three
      votes to one. Hence \(z\succeq^M x\succeq^M y\), but
      \(z\not\succeq^M y\): the weak majority relation is not
      transitive. Its weak winners are \(x,y\), its central peaks
      are \(x,y\), and it has no strict Condorcet winner.
      The annex also constructs this same failure on the full interval
      \([0,1]\) with strict single-peaked individual preferences.
\end{enumerate}
\end{hbtheorem}

\section{Compatibility with Arrow's theorem}
\label{sec:p17-single-peaked-arrow}

\begin{hbtheorem}{thm:p17-single-peaked-arrow}{proved=annex:p17.single-peaked-arrow}{A nondictatorial rule on a proper restricted domain}
On \(X=[0,1]\) with an odd electorate \(n\geq3\), the function
\(F_{\mathrm{SP}}\) returns complete transitive social preferences,
follows every unanimous strict comparison, and is anonymous under
permutations of voter identities. It satisfies IIA when that condition
is quantified over pairs of admissible single-peaked profiles. It is
not dictatorial on this restricted domain.

The domain \(\mathcal D_{\mathrm{SP}}\) is a proper subset of
\(\mathcal L([0,1])^N\). Consequently these properties are compatible
with Theorem~\ref{thm:p17-social-arrow}, which assumes an unrestricted
strict-profile domain. The restriction changes which profiles the rule
must handle; it does not remove transitivity or change the definition of IIA.
For \(n=1\), the rule agrees with its sole voter's ordering, as expected.
\end{hbtheorem}

\begin{proofboundary}{def:p17-single-peaked-next}{Further preference domains and strategic questions}
This chapter proves the source's odd-electorate result for strict
single-peaked preferences on the full interval, its stated ordered-set
extension, the median conclusions and the supplied-axis finite method.
Allowing individual indifference requires a newly specified domain and
proof; the earlier binary May theorem's weak-input convention is not
silently carried into this theorem. Finding a common axis, other domain
restrictions, multidimensional alternatives, strategic misreporting and
the source's remaining social-choice problems require their own treatments.
A Condorcet-winner theorem is not by itself a theorem that reporting the
true preference is a best strategy in a specified mechanism.
\end{proofboundary}
